%! AP Statistics — Unit 8, Lesson 8.3 — Exit Ticket (Answer Key)
\documentclass[10pt]{article}
\usepackage{apstats-article}
\usepackage{apstats-key}

\begin{document}

\pageheader{Unit 8, Lesson 8.3}{Exit Ticket --- Answer Key}
\namedateperiod

\vspace{0.1in}

A bag of 160 colored candies claims that each of its 4 colors makes up exactly 25\% of
the candies. A student counts the colors and records:

\renewcommand{\arraystretch}{1.4}
\begin{center}
\begin{tabular}{lcccc}
    \textbf{Color}            & Red & Blue & Green & Yellow \\
    \hline
    \textbf{Observed ($O$)}   & 46  & 38   & 42    & 34     \\
    \textbf{Expected ($E$)}   & 40  & 40   & 40    & 40     \\
    $\boldsymbol{(O-E)^2/E}$  & 0.9 & 0.1  & 0.1   & 0.9    \\
\end{tabular}
\end{center}

\vspace{0.05in}
$\chi^2 = 0.9 + 0.1 + 0.1 + 0.9 = 2.0$, \quad $df = 3$, \quad p-value $> 0.10$

\vspace{0.12in}

\begin{enumerate}

  \item Using the chi-square table, the p-value for $df = 3$ and $\chi^2 = 2.0$ is
        greater than 0.10. \textbf{Interpret this p-value in context.}

        \ansline{The probability of obtaining a chi-square statistic of 2.0 or larger, given that the null hypothesis (each color makes up exactly 25\% of all candies in the bag) is true, is greater than 0.10.}

  \item At $\alpha = 0.05$, \textbf{state a complete conclusion} for this chi-square
        goodness-of-fit test.

        \ansline{Because $p > 0.10 > \alpha = 0.05$, we fail to reject $H_0$. We do not have convincing evidence that the color distribution differs from 25\% for each color.}

\end{enumerate}

\end{document}
